% Pursuing Stacks (cotes 134-2 to 134-8): archive cote 134-2, site page 60 (PDF page 61).
% appendix (letters to Breen). Draft compared: ps1.tex:1431-1447.
% AI-transcribed by Claude Opus 5.5 (claude-opus-5-5, adaptive thinking, effort high), Message Batches API, batch ps1 (msgbatch_01Rv5hCaNTtUt4JHx1yqoHPL), 2026-11-12; not yet checked by eye.
\scanpage{134-2/60}
\mspage{- 45 -}

(4) This works \inserted{also,} for $f : X \to S$ finitely presented separated if $f$ has the
\margin{espacer} properties of ``cohomological local triviality'' \inserted{(properties ``local upstairs'')} (\deleted{enhaut local properties}) for
example $f$ \textit{smooth} ; one finds that $H\hand{_*}\hand{(}X/S\hand{)} = Rf\hand{_!}f^{\hand{!}}(\Lambda_s)$ . \margin{$*$ $*$ $*$ (three hand stars below the line, under ``that'', ``H'' and ``='')}

(B) Artin-Mazur have studied in a spirit close to yours the \textit{autodualit\over{ies}{y}} of
the Jacobian\deleted{s} of a relative curve $X/S$ . It is necessary to ask them for
precise results, perhaps it works say if $X/S$ is proper and flat of relative
dimension 1 - in any case it is OK on a discrete valuation ring with smooth
\textit{generic} fibre. The special fibre\deleted{s} could be very wild. (I have used their
results in SGA7 to prove, in the case of Jacobians, a duality conjecture on
the group of connected components associated to the Neron models of abelian
varieties dual one to another ...)

Towards the end \deleted{of decade} of the 50's (beginning of 60's?), when the
grand cohomological stuff ($f\hand{_!}$, $f^{\hand{!}}$ , étale cohomology, etc.)\deleted{,} just came out
from darkness, the course given by Serre on the theory of Rosenlicht and Lang
on generalised jacobians and the geometric class field theory (see Serre's book)
and later the ``geometric'' theory of \textit{local} class field theory making use of
pro-algebraic groups (see his article on this subject), made me reflect on
the cohomological formulations of these and other results, which should be
of a ``geometric'' nature, such \deleted{as} \inserted{that} the ``arithmetic'' results \deleted{on the} \inserted{over an} arbitrary
base field (or residual field) $k$ (finite, for example) follow immediately
by descent from the ``geometric'' case of base field $\bar{k}$ . I exchanged letters
with Serre - I don't know if I can find copies - but I recall that \deleted{he suggested} \inserted{I sketched}
projects for some ambitious enough theories on generalised residue\hand{s}\deleted{fiel}ds,
generalised local jacobians, etc., in at least three different directions.
\margin{espacer} But I have never, in spite of numerous attempts, \deleted{come to a satisfactory -} \inserted{succeeded in mobilizing}
\deleted{position with regard to} \inserted{someone for developing one of} these programmes. Here \deleted{are several} \inserted{a few} words on them:

\deleted{(\unsure{App.11|App.II})} \inserted{(C) In the situation} \deleted{In the situation} where $X$ is of finite type \deleted{on a} \inserted{over a} \textit{field} $k$ , construction \margin{$*$ $*$ $*$ (three hand stars above the line)}
of a complex of generalised jacobians $J_{*X/k}$ (of length equal to dim $X$ ).
This is a complex of affine \inserted{commutative pro-}algebraic \deleted{pro}groups on $k$ , with the exception of

\margin{60}
